\chapter{Desarrollo e Implementación de los Modelos}[Desarrollo e Implementación de los Modelos]
\label{cp:modelos}

Una vez consolidada la base de datos dual y garantizada la coherencia geométrica píxel a píxel de las muestras, el foco de este trabajo 
se centra en el diseño, configuración y desarrollo de las arquitecturas de aprendizaje profundo. En el ámbito de la teledetección, la 
elección y el diseño de la estructura de una red neuronal no solo dependen de la naturaleza de los datos de entrada, sino del objetivo 
deseado.


El propósito principal de este capítulo es desglosar las decisiones técnicas y de ingeniería de software adoptadas para materializar las 
herramientas predictivas de este proyecto. El flujo de trabajo se ha estructurado de manera progresiva, comenzando por la configuración 
del entorno de computación. Posteriormente, el diseño experimental se divide en dos aproximaciones metodológicas diferenciadas y complementarias:
un modelo inicial de clasificación global (por parche) y un modelo de segmentación semántica (por píxel).


El desarrollo del primer modelo (un clasificador a nivel de parche) es una aproximación para intuir si las \acrshort{cnn} son capaces de extraer
los patrones macroscópicos y firmas espectrales conjuntas a partir de cuatro bandas seleccionadas de Sentinel-2. Mientras que el segundo, se aborda
una arquitectura basada en una U-net orientada a la segmentación semántica píxel a píxel. Este segundo modelo pretende aumentar la resolución de 
la clasificación, pasando de la clasificación global a la clasificación a nivel de píxel. A lo largo de las siguientes secciones se detallarán
la carga de los datos, la arquitectura de ambas redes, las restricciones que provocaron una redimensión de los parches y las estrategias 
a la hora de controlar las funciones de pérdida, buscando suavizar el desbalanceo de las clases a nivel nacional. 